%=====================================================================
%  Cross-Modal Correspondence in Continual Spiking Networks:
%  An Engineered Bridge, Not an Emergent Geometry
%  A CONTROLLED-NEGATIVE paper. Every number below is pulled from a
%  results file in this repo; gaps are marked \todo{verify: ...}.
%  Compiles out-of-the-box on Overleaf (pdfLaTeX or LaTeX->DVI->PS->PDF).
%
%  NUMBER PROVENANCE (file -> number), for the author:
%   - results/hypergraph_anchor_v2/report.txt  -> in-domain V2A 84.43/79.97%,
%     held-out 18.75/21.88%, gaps; n=8 seeds.
%   - results/hypergraph_anchor_v2/procrustes.json -> Procrustes 0.0/0.0543/-0.0020/0.0563.
%   - results/structural_diagnostic/rdm_analysis.txt -> r=0.0513 (p=0.7379),
%     vision mag rho=0.1128 (p=0.4564), audio mag rho=-0.0395 (p=0.7953).
%   - results/joint_diagnostic/summary.csv + cell_*.csv -> the 2x2
%     (per-sample V2A/A2V/held-out-4/perm-null/det-gap/probe), seeds 42/123/456
%     (separate) and 42 (shared). n_vis_heldout=3969, n_aud_heldout=909.
%   - results/hypergraph_anchor_v2/run.log (lines ~1294-1298) -> baseline
%     table (Retention/V2A/Held-out per condition).
%   - results/cl_confirmation/cl_confirm.log -> N-MNIST 97.21% (10-class),
%     T1 retention ~97.3% after T2/T3.
%   - results/hypergraph_anchor/run2.log + cl_confirm.log -> SHD 79.64% (20-class).
%=====================================================================
\documentclass[conference]{IEEEtran}

\usepackage{cite}
\usepackage{amsmath,amssymb,amsfonts}
\usepackage{graphicx}
\usepackage{booktabs}
\usepackage{url}
\usepackage{xcolor}
\usepackage[hidelinks]{hyperref}

% Inline TODO marker (clearly visible; no extra package needed).
\newcommand{\todo}[1]{\textcolor{red}{[TODO: #1]}}

\title{Cross-Modal Correspondence in Continual Spiking Networks\\
Is Engineered, Not Emergent:\\
A Controlled Negative Result}

\author{
\IEEEauthorblockN{\todo{Author names}}
\IEEEauthorblockA{\todo{Affiliation, email}}
}

\begin{document}
\maketitle

\begin{abstract}
A widespread intuition---the ``Platonic representation'' view---holds that
models trained on different modalities converge toward a shared underlying
geometry, which would let a system transfer a concept learned for one sense
(e.g.\ a visual digit) to another (e.g.\ an auditory digit) without further
supervision. We test this directly in a neuromorphic setting with two
spiking backbones trained separately on N-MNIST (visual digits) and SHD
(audio digits). We ask a single, sharp question: \emph{does cross-modal
correspondence generalize to digit classes that were never paired across
modalities?} Our answer is a controlled negative. (i) Lightweight
correspondence heads that supervise each modality onto fixed orthonormal
anchors reach $80$--$84\%$ in-domain cross-modal retrieval---but this is
\emph{high by construction}, since both modalities are pushed onto the same
fixed targets; it is not a finding. (ii) On held-out digits $\{6,7,8,9\}$
that never receive a shared target, retrieval is at chance: $18.8\%$ (MLP)
and $21.9\%$ (linear) versus a $25\%$ four-class chance, and a
parameter-free orthogonal Procrustes map retrieves $0.0\%$. (iii) A
representational-dissimilarity (RDM) analysis finds no shared geometry
between the two backbones (vision--audio $r=0.05$, $p=0.74$) and no
magnitude/ordinal axis (vision $\rho=0.11$, audio $\rho=-0.04$). (iv) A
joint-encoder $2{\times}2$ ablation then shows that any residual held-out
alignment is an \emph{architectural} artifact of weight sharing, not of the
learning objective: a shared classifier produces $95\%$ held-out retrieval
with zero cross-modal supervision (a supervision leak, $n{=}1$), while a
shared encoder core alone is seed-unstable ($46/36/0.6\%$ across three
seeds, one collapsing below chance), and---critically---the actual
cross-modal contrastive loss \emph{hurts} held-out retrieval
($\text{OFF}\ge\text{ON}$ in all three seeds). We frame the result through
the zero-shot-learning necessary condition \cite{palatucci2009}: arbitrarily
paired classes cannot transfer without supplied side information, which we
withheld for $\{6,7,8,9\}$ by design.
\end{abstract}

\begin{IEEEkeywords}
spiking neural networks, neuromorphic computing, cross-modal alignment,
zero-shot generalization, representational convergence, controlled negative
\end{IEEEkeywords}

%=====================================================================
\section{Introduction}
%=====================================================================
Neuromorphic spiking neural networks (SNNs) are attractive for edge
deployment because they encode information in sparse, event-driven spikes,
but making them learn \emph{continually} across \emph{multiple senses}
remains hard. A natural ambition for such a system is cross-modal transfer:
having learned that a visual ``3'' and an auditory ``3'' co-occur, can the
network infer---for a class it has \emph{never} seen paired---that a visual
``7'' corresponds to an auditory ``7''? If separately trained representations
truly converge to a shared geometry, as the Platonic representation
hypothesis \cite{platonic} and relative-representation stitching
\cite{relreps} suggest they often do in deep networks, the answer should be
yes: a single map learned on a few classes should extrapolate to the rest.

This paper tests that prediction in a controlled neuromorphic setting and
reports a negative result, dissected mechanistically. We train per-modality
SNN backbones on visual (N-MNIST \cite{nmnist}) and audio (SHD \cite{shd})
digits, build a correspondence interface between them with lightweight
heads supervised onto fixed orthonormal anchors, and then ask whether the
interface---or the backbones themselves---contains a geometry generalizable
to the four held-out digits $\{6,7,8,9\}$ that are never given a shared
target. They do not. Held-out retrieval sits at chance across nonlinear,
linear, and parameter-free maps; the two backbones share no relational
structure (RDM $r=0.05$); and a joint-encoder ablation shows that the only
thing that produces apparent held-out alignment is shared \emph{weights},
not the cross-modal learning signal.

We are deliberate about framing. The in-domain correspondence accuracy
($80$--$84\%$) is \emph{high by construction}: both modalities are
supervised against the same fixed anchors, so they land near each other by
design. We never present it as a discovery. Our contribution is the rigorous
\emph{negative} and its mechanism.

\noindent\textbf{Contributions.}
\begin{itemize}
  \item A controlled held-out protocol (train correspondence on digits
        0--5, test on 6--9) with a deterministic null, three map families,
        and $n{=}8$ seeds, showing held-out retrieval at chance and Procrustes
        $0\%$---correspondence is engineered, not a discovered shared geometry.
  \item An RDM root-cause analysis: the separately trained backbones share
        no relational structure ($r=0.05$, $p=0.74$) and carry no
        ordinal/magnitude axis, explaining why no map extrapolates.
  \item A joint-encoder $2{\times}2$ (shared/separate classifier $\times$
        contrastive loss on/off) that isolates the locus of any residual
        alignment: it is architectural (a shared-classifier supervision leak;
        an unstable shared core), and the contrastive loss actively
        \emph{hurts} held-out generalization ($\text{OFF}\ge\text{ON}$ in
        $3/3$ seeds).
  \item An interpretation via the zero-shot-learning necessary condition
        \cite{palatucci2009}: cross-modal transfer of arbitrarily paired
        classes requires side information, which we deliberately withheld.
\end{itemize}

%=====================================================================
\section{Related Work}\label{sec:related}
%=====================================================================
\textbf{Representational convergence.} The Platonic representation
hypothesis \cite{platonic} argues that models across modalities and
objectives trend toward a shared representation of reality; relative
representations \cite{relreps} exploit this by stitching separately trained
latent spaces with a single learned or orthogonal map, enabling zero-shot
model merging. These results are established in dense deep networks at
scale; whether they extend to small, separately trained \emph{spiking}
backbones on heterogeneous neuromorphic sensors is the question we
operationalize. Our negative result is a controlled counterpoint:
convergence is not automatic here, and a single rotation does not
extrapolate to unseen classes.

\textbf{Zero-shot learning.} Zero-shot transfer to unseen classes requires
that those classes be representable through shared side information
\cite{palatucci2009}---attributes, embeddings, or (as in our anchors) fixed
target vectors. \todo{xian: cite ZSL survey (Xian et al.) here and verify
bibliography entry.} Our held-out classes are deliberately given no shared
target; chance-level cross-modal retrieval is therefore the
\emph{predicted} outcome, not an algorithmic failure, and we report it as
such.

\textbf{Cross-modal alignment and neuromorphic memory.} Contrastive
objectives align modalities in a shared space (we use the supervised
contrastive form \cite{supcon}); prior neuromorphic work studies
modality-dependent memory in cross-modal SNNs \cite{paper1}. Our joint
encoder reuses a shared core and contrastive loss, but we find the
contrastive signal does not generalize to unpaired classes.

%=====================================================================
\section{Methods}\label{sec:methods}
%=====================================================================
\subsection{Datasets and backbones}
We use N-MNIST \cite{nmnist} (visual event frames, $34{\times}34{\times}2$,
densified to $25$ time steps) and SHD \cite{shd} (audio, densified to
$100{\times}700$, labels reduced modulo $10$). Per-modality SNN backbones
use leaky integrate-and-fire (LIF) neurons with leak $\beta{=}0.9$ and
rate-coded readouts (spike counts summed over time); the visual stem is
convolutional and the audio stem is fully connected, both emitting
$512$-dimensional features. Standalone single-task accuracies are
$97.21\%$ for N-MNIST (10-class head) and $79.64\%$ for SHD; the SHD number
reflects a $20$-class output head, not lower feature quality, and is stated
only for reference.

\subsection{Frozen-backbone correspondence with fixed anchors}
\label{sec:method:corr}
We freeze both backbones and train lightweight correspondence heads that
project the $512$-d features into a $64$-d anchor space: a linear head
\texttt{Linear(512,64)} and an MLP head
\texttt{Linear(512,256)$\to$GELU$\to$Linear(256,64)}. Ten fixed orthonormal
anchor vectors (one per digit, from a QR decomposition) serve as
\emph{immutable} targets. Each modality's L2-normalized projection $z$ is
trained to its digit anchor with a cosine-classifier cross-entropy loss,
\begin{equation}
\mathcal{L}_{\mathrm{corr}} \;=\; \mathrm{CE}\!\left( \frac{z A^{\!\top}}{\tau},\; y \right),
\quad z = \tfrac{\mathrm{proj}(\mathrm{backbone}(x))}{\|\mathrm{proj}(\mathrm{backbone}(x))\|},
\label{eq:loss}
\end{equation}
with temperature $\tau{=}0.1$ and fixed anchor matrix $A$. Because the
targets are fixed (not learned), they cannot drift---the failure mode of a
shared, trainable target---and because nothing on the inference path is
updated, retention of the original tasks is zero-forgetting by construction
(verified at $97.3\%$ retention of N-MNIST through the task sequence).

\subsection{Held-out protocol}
Digits are split into train classes $\{0,\dots,5\}$ and held-out classes
$\{6,7,8,9\}$. Correspondence is trained on $0$--$5$ only; $\{6,\dots,9\}$
receive \emph{no shared target} and are reserved to test generalization. We
evaluate three map families: the nonlinear MLP head, the linear head, and a
parameter-free orthogonal Procrustes rotation $R$ fit on train-class
prototypes and applied to held-out prototypes. We also compute a
deterministic null (mean cosine to non-matching prototypes, no shuffle) and
report the gap (same$-$null).

\subsection{Joint-encoder $2{\times}2$ ablation}
\label{sec:method:joint}
To separate architectural sharing from the learning objective, we train a
single encoder from scratch on both modalities (visual/audio stems into a
\emph{shared} core \texttt{Linear(256,128)$\to$GELU}, L2-normalized) with
within-modality digit classification on all $0$--$9$. We cross two axes:
classifier \{separate per modality vs.\ one shared head\} and cross-modal
contrastive loss \{on (supervised contrastive \cite{supcon}, temperature
$0.1$, applied to $0$--$5$ only) vs.\ off\}. The contrastive-off row is
decisive: if held-out $6$--$9$ align with no cross-modal supervision, the
alignment is architectural.

\subsection{Metrics and statistical controls}
Held-out retrieval is computed \emph{per sample}: each held-out test sample
is matched to its nearest opposite-modality class prototype, with chance
$=10\%$ over all ten classes (we also report the four-class variant at
$25\%$ chance for cross-reference). We use a permutation null (relabel
prototypes, $200$ iterations) as an empirical chance estimate and a
binomial test against the analytical $10\%$ level ($n_{\mathrm{vis}}{=}3969$,
$n_{\mathrm{aud}}{=}909$ held-out samples). An untrained-encoder control and
the deterministic null guard against metric artifacts. We pre-registered
the expectation that held-out retrieval be at chance and report deviations
honestly.

%=====================================================================
\section{Results}\label{sec:results}
%=====================================================================
\subsection{In-domain correspondence is high---by construction}
The correspondence heads reach $84.43\%\pm1.17$ (MLP) and $79.97\%\pm0.99$
(linear) in-domain visual$\to$audio retrieval over $n{=}8$ seeds, at
$97.3\%$ retention. \todo{verify: in-domain A2V and prototype-alignment
numbers are not in the results files; omit or fill from a sourced file.}
We stress that these numbers are high \emph{by construction}: both
modalities are supervised onto the \emph{same} fixed anchors, so their
projections necessarily land near each other. This is the expected behavior
of the training objective, not evidence of a discovered shared geometry, and
we do not present it as a finding.

\subsection{Held-out generalization fails across all map families}
On digits $\{6,\dots,9\}$, never given a shared target, retrieval collapses
to chance (Table~\ref{tab:heldout}): $18.75\%\pm16.54$ (MLP) and
$21.88\%\pm19.52$ (linear) against a $25\%$ four-class chance, with gaps
indistinguishable from zero. The parameter-free orthogonal Procrustes map
retrieves $0.0\%$ (same-class alignment $0.054$ vs.\ null $-0.002$). Because
the linear map---the least constrained of the three---also fails, the
deficit is not an under-capacity artifact; because Procrustes fails, no
global rotation relates the two held-out geometries.

\begin{table}[t]
\centering
\caption{Held-out generalization (digits $6$--$9$, never supervised; four-class
retrieval chance $=25\%$). All map families sit at chance; Procrustes
retrieves $0\%$. Mean $\pm$ std over $n{=}8$ seeds. Sources:
\texttt{results/hypergraph\_anchor\_v2/report.txt}, \texttt{procrustes.json}.}
\label{tab:heldout}
\begin{tabular}{lccc}
\toprule
Map & Held-out retrieval & Gap (same$-$null) & Verdict \\
\midrule
MLP head (nonlinear)    & $18.75 \pm 16.54\%$ & $-0.044 \pm 0.147$ & at chance \\
Linear head             & $21.88 \pm 19.52\%$ & $\phantom{-}0.039 \pm 0.084$ & at chance \\
Procrustes (orthogonal) & $0.0\%$             & $\phantom{-}0.056$           & at chance \\
\bottomrule
\end{tabular}
\end{table}

\subsection{Root cause: no shared geometry (RDM)}
A representational-dissimilarity analysis explains why nothing extrapolates:
the two backbones do not share relational structure. The vision--audio RDM
correlation is $r=0.0513$ ($p=0.7379$); nor is there a shared ordinal/
magnitude axis---Spearman $\rho$ between embedding magnitude and digit value
is $0.1128$ ($p=0.4564$) for vision and $-0.0395$ ($p=0.7953$) for audio
(source: \texttt{results/structural\_diagnostic/rdm\_analysis.txt}). With
neither a shared global geometry nor a latent ordering to extrapolate along,
no map fit on $\{0,\dots,5\}$ can reach $\{6,\dots,9\}$.

\subsection{The joint-encoder $2{\times}2$: alignment is architectural}
If frozen backbones share no geometry, does \emph{co-training} a single
encoder on both modalities create one? Table~\ref{tab:2x2} reports the
per-sample held-out metric (chance $10\%$). Three observations:

\textbf{(a) A shared classifier is a supervision leak.} With one shared
$0$--$9$ head and \emph{no} cross-modal loss, held-out retrieval is
$95.4\%$ (V2A) / $88.7\%$ (A2V). This is not learning: the shared head is
trained to output ``7'' in both modalities, which forces visual-$7$ and
audio-$7$ into the same slot. It is the artifact that \emph{explains}
apparent generalization, not a positive result, and it is $n{=}1$.

\textbf{(b) The shared core alone is seed-unstable.} With \emph{separate}
classifiers and no cross-modal loss, retrieval is $46.0\%$, $35.9\%$, and
$0.6\%$ across seeds $42/123/456$ (mean $27.5\%\pm19.5$). Two seeds show a
moderate shared-core effect; the third collapses \emph{below} chance
(binomial $z\approx-20$ vs.\ $10\%$). This is not a robust phenomenon.

\textbf{(c) The contrastive loss hurts.} Turning the cross-modal
contrastive loss \emph{on} does not improve held-out retrieval---it reduces
it: $\text{OFF}\ge\text{ON}$ in all three seeds ($46{\ge}20$,
$36{\ge}26$, $0.6{\ge}0.4$; means $27.5\%$ vs.\ $15.2\%$). The genuine
cross-modal learning signal over-specializes the shared geometry to the
trained classes $\{0,\dots,5\}$ and does not transfer to $\{6,\dots,9\}$.

\begin{table}[t]
\centering
\caption{Joint-encoder $2{\times}2$: per-sample held-out retrieval
(digits $6$--$9$; chance $=10\%$ over all ten classes; permutation null
$\approx10\%$). The only cell well above chance is the shared-classifier
row---a supervision leak, not learning. Sources:
\texttt{results/joint\_diagnostic/summary.csv}, \texttt{cell\_*.csv}.}
\label{tab:2x2}
\begin{tabular}{llcccc}
\toprule
Classifier & SupCon & V2A (\%) & A2V (\%) & det-gap & seeds \\
\midrule
Separate & off & $27.5 \pm 19.5$ & $42.6 \pm 14.9$ & $0.084$ & 3 \\
Separate & on  & $15.2 \pm 10.7$ & $14.3 \pm 10.6$ & $-0.021$ & 3 \\
Shared   & off & $95.4$          & $88.7$          & $1.085$ & 1 \\
Shared   & on  & $94.8$          & $86.1$          & $1.092$ & 1 \\
\midrule
\multicolumn{6}{l}{Permutation null (all cells): $\approx9.6$--$10.1\%$.} \\
\multicolumn{6}{l}{Within-modality $6$--$9$ probe (gate): vis $\sim94$--$95\%$, aud $\sim86$--$89\%$ (passes).} \\
\multicolumn{6}{l}{Separate-row per-seed V2A: off $\{46.0, 35.9, 0.6\}$; on $\{19.7, 25.5, 0.4\}$.} \\
\bottomrule
\end{tabular}
\end{table}

\subsection{Baseline context}
For reference, the correspondence interface achieves $84.4\%$/$80.0\%$
in-domain V2A (MLP/linear) at $97.3\%$ retention, against
no-alignment ($\sim13.3\%$ V2A) and frozen-backbone ($14.8\%$) baselines
(source: \texttt{results/hypergraph\_anchor\_v2/run.log}). The in-domain
gap is an effect of the anchors, not of representational convergence; the
held-out numbers above are the load-bearing comparison.

% Optional figures (drop the PDFs in and uncomment).
% \begin{figure}[t]
% \centering
% % \includegraphics[width=\columnwidth]{fig_rdm.pdf}
% \fbox{\parbox{0.9\columnwidth}{\centering\vspace{6pt}RDM vision vs.\ audio\\(drop \texttt{fig\_rdm.pdf})\vspace{6pt}}}
% \caption{Vision--audio representational dissimilarity is uncorrelated ($r=0.05$). Source: \texttt{results/structural\_diagnostic/rdm\_*.npy}.}
% \label{fig:rdm}
% \end{figure}
%
% \begin{figure}[t]
% \centering
% % \includegraphics[width=\columnwidth]{fig_2x2.pdf}
% \fbox{\parbox{0.9\columnwidth}{\centering\vspace{6pt}2$\times$2 held-out retrieval\\(drop \texttt{fig\_2x2.pdf})\vspace{6pt}}}
% \caption{Joint-encoder $2{\times}2$: only the shared-classifier cell clears chance, and it is a leak. Source: \texttt{results/joint\_diagnostic/summary.csv}.}
% \label{fig:2x2}
% \end{figure}

%=====================================================================
\section{Discussion}\label{sec:disc}
%=====================================================================
The locus of any cross-modal alignment we observe is \emph{architectural
weight sharing}, not the learning objective. A shared output head
manufactures alignment by supervising both modalities to the same label
slot; a shared encoder core produces a weaker, seed-unstable version; and
the one signal that would constitute genuine cross-modal learning---the
contrastive loss---does not transfer to held-out classes and in fact harms
them. Read together with the frozen-backbone result (Procrustes $0\%$, RDM
$r=0.05$), the conclusion is that separately trained spiking backbones on
these modalities do not converge to a generalizable cross-modal geometry,
and co-training alone does not create one.

This is consistent with the zero-shot-learning necessary condition
\cite{palatucci2009}: zero-shot transfer to unseen classes requires shared
side information tying the unseen classes to the seen ones. Our fixed
orthonormal anchors supply such side information for the trained classes
$\{0,\dots,5\}$ (which is exactly why in-domain retrieval is high), but we
deliberately withhold it for $\{6,\dots,9\}$. Chance-level held-out
retrieval is therefore the theoretically predicted outcome. The
contrastive loss over-specializing to trained classes is the
empirical signature of this: with nothing tying $6$--$9$ to $0$--$5$,
optimizing the shared geometry for $0$--$5$ moves it \emph{away} from a
configuration that would incidentally also serve $6$--$9$.

We are explicit about what these results are \emph{not}. They are not a
discovery of cross-modal alignment; they are not a claim that convergence
is impossible in general---only that it is not automatic here, and that
apparent transfer can be fully explained by architecture and supervision
leakage. The in-domain numbers are high by construction and are reported
only to characterize the interface.

%=====================================================================
\section{Conclusion and Limitations}\label{sec:conc}
%=====================================================================
We presented a controlled negative result: cross-modal correspondence
between separately trained neuromorphic SNN backbones does not generalize
to held-out digit classes, across nonlinear, linear, and parameter-free
maps; the backbones share no relational geometry; and the only source of
apparent held-out alignment in a joint encoder is architectural weight
sharing (a shared-classifier leak and an unstable shared core), while the
cross-modal contrastive objective actively hurts held-out transfer. The
result is exactly what the zero-shot necessary condition predicts when
side information is withheld.

\noindent\textbf{Limitations.}
The shared-classifier leak ($95.4\%$) is reported at $n{=}1$ seed and,
though extreme, should be confirmed over more seeds; the shared-core effect
is demonstrably seed-unstable (one of three seeds collapses below chance);
correspondence is evaluated on digits only; and the contrastive-loss
result holds for one temperature and capacity. The honest route to
\emph{genuine} discovery is a compositional/shared side-information axis
(e.g.\ a supplied magnitude or attribute) tying unseen classes to seen
ones---future work we name but do not run here.

%=====================================================================
\section*{Acknowledgments}
\todo{acknowledgments / funding if any.}

%=====================================================================
%  REFERENCES
%  Real citation keys reused from paper_skeleton.tex; details flagged
%  \todo{verify} where author/venue/year must be confirmed before submission.
%  Palatucci2009 and Xian are NEW (not in the skeleton) and must be verified.
%=====================================================================
\begin{thebibliography}{99}
\bibitem{paper1}
\todo{Paper 1 --- Modality-Dependent Memory Mechanisms in Cross-Modal
Neuromorphic Computing (self-cite). Verify title/authors/venue.}

\bibitem{platonic}
M.~Huh, B.~Cheung, P.~Wang, and P.~Isola, ``The Platonic Representation
Hypothesis,'' \emph{arXiv:2405.07987}, 2024. \todo{verify}

\bibitem{relreps}
L.~Moschella, V.~Maiorca, M.~Fumero, A.~Norelli, F.~Locatello, and
E.~Rodol{\`a}, ``Relative Representations Enable Zero-Shot Latent Space
Communication,'' in \emph{ICLR}, 2023. \todo{verify}

\bibitem{palatucci2009}
M.~Palatucci, G.~Hinton, T.~Mitchell, and R.~Salakhutdinov, ``Zero-shot
Learning with Semantic Output Codes,'' in \emph{Advances in Neural
Information Processing Systems (NeurIPS)}, 2009.
\todo{verify --- NEW key, not in skeleton.}

\bibitem{xian}
\todo{Xian \emph{et al.}, zero-shot learning survey (cite in
Sec.~\ref{sec:related}). Provide authors, title, venue, year.}

\bibitem{supcon}
P.~Khosla, P.~Teterwak, C.~Wang, A.~Sarna, Y.~Tian, P.~Isola, A.~Maschinot,
C.~Liu, and D.~Krishnan, ``Supervised Contrastive Learning,'' in
\emph{NeurIPS}, 2020. \todo{verify}

\bibitem{shd}
B.~Cramer, D.~Strotdreses, S.~Senn, and W.~Maass, ``The Heidelberg Spiking
Data Sets for the Systematic Evaluation of Spiking Neural Networks,''
\emph{IEEE Trans.\ Neural Netw.\ Learn.\ Syst.}, 2020. \todo{verify}

\bibitem{nmnist}
G.~Orchard, A.~Jayawant, G.~K.~Cohen, and N.~Thakor, ``Converting Static
Image Datasets to Spiking Neuromorphic Datasets Using Saccades,''
\emph{Front.\ Neurosci.}, 2015. \todo{verify}
\end{thebibliography}

\end{document}
